% =====================================================================
% PRJ-23 - Aircraft Preliminary Design
% Homework 03 - Perfil Aerodinamico
% Grupo NJ-0502 - Instituto Tecnologico de Aeronautica
% Compilar com: pdflatex main.tex (duas vezes)
% =====================================================================
\documentclass[11pt,a4paper]{article}

\usepackage[utf8]{inputenc}
\usepackage[T1]{fontenc}
\usepackage[brazil]{babel}
\usepackage{lmodern}
\usepackage{amsmath,amssymb}
\usepackage{graphicx}
\usepackage{booktabs}
\usepackage{longtable}
\usepackage{float}
\usepackage{xcolor}
\usepackage{microtype}
\usepackage[margin=2.4cm]{geometry}
\usepackage[font=small,labelfont=bf]{caption}
\usepackage{fancyhdr}
\usepackage{pdflscape}
\usepackage[colorlinks=true,linkcolor=black,citecolor=black,urlcolor=blue!60!black]{hyperref}

\pagestyle{fancy}
\fancyhf{}
\fancyhead[L]{\small PRJ-23 --- Projeto Preliminar de Aeronave}
\fancyhead[R]{\small Lab 03 --- Perfil Aerodinâmico}
\fancyfoot[C]{\thepage}
\renewcommand{\headrulewidth}{0.4pt}

\newcommand{\code}[1]{\texttt{\small #1}}

\begin{document}

\begin{titlepage}
  \centering
  \vspace*{1cm}

  {\Large \textbf{Instituto Tecnologico de Aeronautica}} \\[0.4cm]
  {\large Departamento de Engenharia Aeronautica} \\[1.5cm]

  \rule{\linewidth}{1pt} \\[0.4cm]
  {\LARGE \textbf{PRJ-23 --- Laboratorio 03}} \\[0.3cm]
  {\Large NJ-0502} \\[0.3cm]
  \rule{\linewidth}{1pt} \\[2cm]

  \begin{tabular}{rl}
    \textbf{Integrantes:}  & Bruno Pimentel Santos \\
                           & Douglas Rodrigues Cabral \\
                           & Emanuel Garcia Alarcon \\
                           & Gabriel Souza Bonanni \\
                           & Joao Pedro Cruz Magalhaes \\
                           & Yuri Miranda Peres \\[0.6cm]
    \textbf{Professor:}    & Maj Eng Ney Rafael Secco \\[0.6cm]
    \textbf{Data:}         & 13 de setembro de 2026 \\
  \end{tabular}

  \vfill
  {\small São José dos Campos, 2026}
\end{titlepage}

\newpage

\begin{abstract}
\noindent
Este documento e um modelo de relatorio para o Laboratorio 03 de PRJ-23.
O texto foi estruturado no mesmo estilo adotado nos Labs 01 e 02, com uma
introducao geral, seccoes separadas por etapa do trabalho, espacos para
figuras, tabelas e discussoes, e uma conclusao final. O modelo esta
organizado para receber os resultados das etapas de parametrizacao do
perfil, analise com o XFOIL e analise/otimizacao com o Euler solver.
\end{abstract}

\vspace{0.3cm}
\hrule
\vspace{0.5cm}



% =====================================================================
% Homework 03 --- Parametrizacao e geracao de perfis
% Incluido por main.tex. Tabelas em tex_airfoil_mod/ e figuras em
% resultados_airfoil_mod/ (mesmo nivel que este arquivo e o main.tex).
% Codigos em airfoil_mod/.
% =====================================================================

\section{Ponto de projeto da seção de referência}
\label{sec:ponto-projeto}

O primeiro passo do laboratório é fixar o ponto de projeto em que o
aerofólio será otimizado. A aeronave de referência é a NJ-0502 com os
inputs já atualizados para o ótimo do Lab~02. Adotou-se o
\textbf{meio do cruzeiro}: altitude e Mach de cruzeiro da missão de
projeto, e o peso médio aritmético entre o início e o fim dessa etapa da missão.

\subsection{Altitude, Mach e área de referência}

Na missão de projeto da NJ-0502,
\begin{equation}
h = 35\,000~\mathrm{ft},
\qquad
M_{\infty} = 0{,}85,
\qquad
S_{\mathrm{ref}} = S_w = 386{,}48~\mathrm{m}^2.
\end{equation}
A área de referência é a área alar otimizada no Lab~02.

\subsection{Atmosfera no ponto de projeto}

Com a atmosfera padrão ISA em $h = 35\,000$~ft ($10\,668$~m),
\begin{equation}
\rho_{\infty} = 0{,}3805~\mathrm{kg/m}^3,
\qquad
a_{\infty} = 296{,}59~\mathrm{m/s},
\qquad
\mu_{\infty} = 1{,}445\times 10^{-5}~\mathrm{Pa\,s}.
\end{equation}
A velocidade de voo correspondente é
\begin{equation}
V_{\infty} = M_{\infty}\,a_{\infty}
= 0{,}85\times 296{,}59
= 252{,}10~\mathrm{m/s}.
\end{equation}

\subsection{Peso no meio do cruzeiro}

O \emph{Design Tools} decompõe a missão em frações de massa. Para a
categoria \emph{transport},
\begin{equation}
\begin{aligned}
M_{f,\mathrm{start}} &= 0{,}990,&
M_{f,\mathrm{taxi}} &= 0{,}990,&
M_{f,\mathrm{TO}} &= 0{,}995,&
M_{f,\mathrm{climb}} &= 0{,}980.
\end{aligned}
\end{equation}
Com o MTOW convergido $W_0 = 276\,276$~kgf, o peso no \emph{início} do
cruzeiro é
\begin{equation}
W_{\mathrm{ini}}
= W_0\,M_{f,\mathrm{start}}\,M_{f,\mathrm{taxi}}\,M_{f,\mathrm{TO}}\,M_{f,\mathrm{climb}}
= 264\,036~\mathrm{kgf}.
\end{equation}
A fração de cruzeiro vem da equação de Breguet implementada em
\code{fuel\_weight},
\begin{equation}
M_{f,\mathrm{cruise}} = 0{,}6992,
\end{equation}
de modo que o peso no \emph{fim} do cruzeiro fica
\begin{equation}
W_{\mathrm{fim}} = W_{\mathrm{ini}}\,M_{f,\mathrm{cruise}}
= 184\,607~\mathrm{kgf}.
\end{equation}
O ponto de projeto adotado é a média aritmética entre esses extremos:
\begin{equation}
W
= \frac{W_{\mathrm{ini}} + W_{\mathrm{fim}}}{2}
= 224\,321~\mathrm{kgf}.
\end{equation}

\subsection{Coeficiente de sustentação da aeronave}

Em voo nivelado,
\begin{equation}
C_L
= \frac{2W}{\rho_{\infty}\,S_{\mathrm{ref}}\,V_{\infty}^{2}}
= 0{,}4710.
\end{equation}
Seguindo o enunciado, a seção de referência herda esse valor:
\begin{equation}
c_{l\mathrm{ref}} = C_L = 0{,}4710.
\end{equation}

\subsection{Corda e espessura da seção de referência}

A corda de referência é a corda média aerodinâmica (MAC) da asa
otimizada,
\begin{equation}
c_{\mathrm{ref}} = c_{m,w}
= \frac{2}{3}\,c_{r,w}\,\frac{1+\lambda_w+\lambda_w^{2}}{1+\lambda_w}
= 6{,}918~\mathrm{m},
\end{equation}
com $\lambda_w = 0{,}20$. A espessura relativa mínima para a otimização
é a média entre raiz e ponta da asa da NJ-0502,
\begin{equation}
\left(\frac{t}{c}\right)_{\mathrm{ref}}
= \frac{(t/c)_{r} + (t/c)_{t}}{2}
= \frac{0{,}16 + 0{,}08}{2}
= 0{,}12.
\end{equation}

\subsection{Reynolds da seção de referência}

\begin{equation}
Re_{\mathrm{ref}}
= \frac{\rho_{\infty}\,V_{\infty}\,c_{\mathrm{ref}}}{\mu_{\infty}}
= 4{,}592\times 10^{7}.
\end{equation}

\subsection{Tabela-resumo}

Os dados da seção de referência estão reunidos na Tabela~\ref{tab:ponto}.

\begin{table}[H]
\centering
\caption{Dados para a seção de referência.}
\label{tab:ponto}
\small
\begin{tabular}{lll}
\toprule
Parâmetro & Valor & Descrição \\
\midrule
$h$ & $35\,000$ & Altitude do ponto de projeto [ft] \\
$M_{\infty}$ & $0{,}85$ & Mach do ponto de projeto \\
$W$ & $224\,321{,}2$ & Peso da aeronave no ponto de projeto [kgf] \\
$S_{\mathrm{ref}}$ & $386{,}48$ & Área de referência da aeronave [m$^2$] \\
$\rho_{\infty}$ & $0{,}3805$ & Densidade do ar no ponto de projeto [kg/m$^3$] \\
$a_{\infty}$ & $296{,}59$ & Velocidade do som no ponto de projeto [m/s] \\
$c_{\mathrm{ref}}$ & $6{,}918$ & Corda da seção de referência [m] \\
$c_{l\mathrm{ref}}$ & $0{,}4710$ & Coeficiente de sustentação da seção de referência \\
$Re_{\mathrm{ref}}$ & $4{,}592 \times 10^{7}$ & Reynolds para a seção de referência \\
$(t/c)_{\mathrm{ref}}$ & $0{,}12$ & Espessura relativa base da seção de referência \\
\bottomrule
\end{tabular}

\end{table}

% ---------------------------------------------------------------------
\section{Parametrização e Geração de Perfis}
\label{sec:airfoil-mod}
% ---------------------------------------------------------------------

\subsection{Objetivo}

Explique aqui qual parametrização foi usada no laboratório
(por exemplo, NACA, CST ou outra), quais variáveis geométricas foram
adotadas e qual era o objetivo desta etapa.

\subsection{Organização dos códigos}

Os códigos desta seção estão em
\code{designTool/lab03\_perfil/airfoil\_mod/}:
\begin{itemize}
\item \code{run\_1\_ponto\_projeto.py} --- ponto de projeto e Tabela~\ref{tab:ponto};
\item demais rotinas de geração/ajuste de perfil.
\end{itemize}

\subsection{Metodologia}

Registre a definição do perfil, os parâmetros geométricos escolhidos,
o procedimento de geração das coordenadas e, se houver, o método de ajuste
do perfil de referência.

\subsection{Resultados}

Insira aqui as figuras e tabelas desta etapa.

\begin{figure}[H]
\centering
% \includegraphics[width=0.85\textwidth]{resultados_airfoil_mod/exemplo.png}
\caption{Substitua esta legenda por uma descrição da figura da etapa de
parametrização.}
\label{fig:airfoil-mod-exemplo}
\end{figure}

\subsection{Discussão}

Discuta a qualidade do ajuste, a interpretação dos parâmetros e eventuais
vantagens ou limitações da parametrização usada.

% =====================================================================
% Homework 03 --- Analise do perfil com XFOIL
% Incluido por main.tex. Tabelas em tex_analise_xfoil/ e figuras em
% resultados_analise_xfoil/ (mesmo nivel que este arquivo e o main.tex).
% Codigos em analise_xfoil/.
% =====================================================================

\section{Analise Aerodinamica com XFOIL}
\label{sec:xfoil}

\subsection{Objetivo}

Descreva o que foi calculado com o XFOIL: polar aerodinamica,
curva $C_L(\alpha)$, arrasto, momento, comparacao entre perfis ou qualquer
outra analise pedida no laboratorio.

\subsection{Configuracao}

Registre aqui as condicoes de escoamento e a forma como a analise foi
conduzida. Exemplos:
\begin{itemize}
\item numero de Reynolds;
\item numero de Mach;
\item faixa de angulo de ataque;
\item criterio de convergencia;
\item configuracao de transicao, se aplicavel.
\end{itemize}

\begin{table}[H]
\centering
\caption{Parametros da analise com o XFOIL.}
\label{tab:xfoil-config}
\begin{tabular}{ll}
\toprule
Parametro & Valor \\
\midrule
Reynolds & preencher \\
Mach & preencher \\
Faixa de $\alpha$ & preencher \\
Observacoes & preencher \\
\bottomrule
\end{tabular}
\end{table}

\subsection{Resultados}

Apresente aqui os graficos e tabelas obtidos nas corridas do XFOIL.

\begin{figure}[H]
\centering
% \includegraphics[width=0.82\textwidth]{resultados_analise_xfoil/polar.png}
\caption{Substitua esta legenda por uma descricao do principal resultado
do XFOIL.}
\label{fig:xfoil-polar}
\end{figure}

\subsection{Discussao}

Comente o comportamento do perfil, a faixa linear de sustentacao,
o arrasto, o ponto de perda, a sensibilidade a Reynolds/Mach e quaisquer
diferencas em relacao a modelos mais simples ou a resultados esperados.

% =====================================================================
% Homework 03 --- Analise e otimizacao com Euler solver
% Incluido por main.tex. Tabelas em tex_eulerblock/ e figuras em
% resultados_eulerblock/ (mesmo nivel que este arquivo e o main.tex).
% Codigos em eulerblock/.
% =====================================================================

\section{Analise e Otimizacao com Euler Solver}
\label{sec:euler}

\subsection{Objetivo}

Explique nesta secao o papel do solver de Euler no laboratorio:
analise do perfil em escoamento inviscido, comparacao com o XFOIL,
estudo de malha, ou otimizacao do perfil.

\subsection{Modelo Numerico}

Descreva aqui:
\begin{itemize}
\item tipo de solver e ordem numerica;
\item malha adotada;
\item condicoes de contorno;
\item criterio de convergencia;
\item variaveis de projeto e restricoes, se houver otimizacao.
\end{itemize}

\begin{equation}
\min_{\mathbf{x}} \; f(\mathbf{x})
\qquad
\text{sujeito a} \qquad
\mathbf{g}(\mathbf{x}) \ge 0.
\end{equation}

\subsection{Resultados}

Inclua aqui os resultados principais desta etapa: historico de convergencia,
campo de pressao, comparacao de perfis, frente de iteracoes, ou historico
da otimizacao.

\begin{figure}[H]
\centering
% \includegraphics[width=0.85\textwidth]{resultados_eulerblock/exemplo.png}
\caption{Substitua esta legenda por uma descricao da principal figura da
etapa com o solver de Euler.}
\label{fig:euler-exemplo}
\end{figure}

\subsection{Discussao}

Comente a convergencia, o custo computacional, a sensibilidade a malha e a
comparacao entre os resultados do Euler solver e os obtidos no XFOIL.




\end{document}
